% TOP_PROBLEM: {"id": 1209, "title": "The Filling Area Conjecture", "category_id": 6, "category": {"id": 6, "name": "geometry", "display_name": "Geometry", "description": "Euclidean and non-Euclidean geometry, geometric structures.", "slug": "geometry", "order_index": 6, "created_at": "2026-07-31T15:26:25.671Z"}, "status": "open"}
% FIRST_POSTED: 2026-09-25
% ENTRY_KIND: statement_only
% !TeX program = lualatex
% Definition and sources only; this file is not an LLM solution attempt.
\documentclass[11pt]{article}
\usepackage[margin=25mm]{geometry}
\usepackage{fontspec}
\setmainfont{Latin Modern Roman}
\usepackage{amsmath,amssymb,mathtools}
\usepackage{unicode-math}
\setmathfont{Latin Modern Math}
\usepackage{xurl,hyperref}
\hypersetup{colorlinks=true,urlcolor=blue}
\setcounter{secnumdepth}{0}
\setlength{\parindent}{0pt}
\setlength{\parskip}{5pt}
\setlength{\emergencystretch}{3em}
\begin{document}
\section*{\texorpdfstring{The Filling Area Conjecture}{The Filling Area Conjecture}}\label{1209}
\subsection{Definitions and mathematical statement}
Let $M$ be a compact orientable Riemannian surface whose boundary is isometric to the circle $C$ of circumference $2\pi$. Write $d_M$ for intrinsic distance allowing paths through $M$, and $d_C$ for the shorter-arc distance along the boundary circle. Suppose no interior path shortens any boundary distance:
\[
 d_M(x,y)=d_C(x,y)\qquad(x,y\in C).
\]
Gromov's filling area assertion is $\operatorname{Area}(M)\ge2\pi$. The unit round hemisphere has this boundary distance and area, fixing the constant. The surface is not required to be a disk: arbitrary orientable topology is allowed. The equality of boundary distances is stronger than prescribing boundary length alone, which would permit very small-area fillings.
\par\smallskip\textit{Formulation and concept references:} \href{https://arxiv.org/abs/2602.17859}{[S1]}.

\subsection{Short English statement}
Must every orientable surface filling a circle without shortening distances between boundary points have at least the area of a round hemisphere?
\subsection{Sources}
\begin{itemize}
\item[S1]Joseph Briggs and Chris Wells, “A Discrete View of Gromov's Filling Area Conjecture,” arXiv:2602.17859 (2026).
\url{https://arxiv.org/abs/2602.17859}.
\textit{Formulation source.}
\end{itemize}
\end{document}
